%------------------------------------------------------------------------------
% 4.4 تصميم قواعد البيانات والتخزين
% المالك: BE (مهمة جانت ٢.٢.٢ — التسليم 31/10/2026)
% الحجم المستهدف: 3-4 صفحات
%
% عناصر مطلوبة:
%   [ ] 4.4.1 قاعدة بيانات ERP (Odoo Education): فلسفة «القراءة فقط عبر
%       العروض الذكية» — لا تعديل لمخطط Odoo إطلاقاً؛ الجداول الأساسية
%       المستفادة (الطلاب، البرامج، الكليات، الرسوم، النتائج...) ومستخدم
%       الربط المخصص بأدنى صلاحية (SELECT على العروض فقط)
%   [ ] 4.4.2 قاعدة بيانات النظام الداخلية + مخطط ER (fig:bayan-er) — TikZ
%       جديد في figures/tikz/bayan-er.tex: المستخدمون، الأدوار، الجلسات،
%       سجل التدقيق، تعريفات العروض الذكية، فهرس المستندات — بالعلاقات
%       والمفاتيح. راجع backend/prisma/schema.prisma في مستودع Bayan للدقة
%   [ ] 4.4.3 مخزن المتجهات Qdrant: المجموعات (Collections)، بنية الحمولة
%       (Payload: المصدر، رقم الصفحة، التصنيف الأمني للوثيقة)، نموذج التضمين
%       multilingual-e5-large وأبعاد المتجهات، واستراتيجية الفهرسة وإعادة
%       البناء عند تحديث المستندات
%
% تحذيرات:
%   - مخطط ER هنا للقاعدة الداخلية الجديدة فقط؛ مخطط Odoo جاهز ومستقر
%     ويُكتفى بجرد الجداول المستخدمة منه دون إعادة رسمه بالكامل
%   - هذا القسم هو موضع «مخططات الكيانات» التي تعد بها مهمة جانت ٢.٢.٤
%------------------------------------------------------------------------------
\section{تصميم قواعد البيانات والتخزين}
\label{sec:db-design}
~

\subsection{قاعدة بيانات \EN{ERP} بيئة التطبيق}
\label{subsec:erp-db}
~

\subsection{قاعدة البيانات الداخلية ومخطط \EN{ER}}
\label{subsec:bayan-db}
~

\subsection{مخزن المتجهات \EN{Qdrant}}
\label{subsec:vector-store}
~
